\section{Sharp conductor reduction and a curve criterion}
\label{sec:sharp-conductor}
All schemes in this article are over $\C$. Put $V_r=H^0(\Pj^1,\mathcal O(r))$,
with $V_r=0$ for $r<0$. For a subspace $U$ of sections, $U^j$ denotes
its multiplication image, not its abstract symmetric power.

\begin{lemma}[The evaluation-kernel bound]\label{lem:sharp-kernel}
Let $U\subset V_n$ be base-point-free, of codimension $c$, and write
$M_U=\ker(U\otimes\mathcal O\to\mathcal O(n))$.
There are integers $a_i\ge1$, with $\sum_i a_i=n$, such that
$M_U=\bigoplus_{i=1}^{n-c}\mathcal O(-a_i)$. If $a=\max_i a_i$, then
for $t\ge0$,
\[
 \operatorname{coker}(U\otimes V_t\to V_{n+t})
       \simeq H^1(M_U(t)),\qquad
 UV_t=V_{n+t}\ \Longleftrightarrow\ t\ge a-1.
\]
In particular $a\le c+1$, so $UV_t=V_{n+t}$ for $t\ge c$.
\end{lemma}
\begin{proof}
The splitting theorem for vector bundles on $\Pj^1$ applies to the
locally free evaluation kernel. Its degree is $-n$ and its rank is
$n-c$. Moreover $H^0(M_U)=0$, because $U\to H^0(\mathcal O(n))$ is
injective. Thus all its summands have negative degree. Twisting the
evaluation sequence by $\mathcal O(t)$ and using $H^1(\mathcal O(t))=0$
gives the asserted cokernel. The formula for $H^1(\mathcal O(t-a_i))$
gives the equivalence. Finally the other $n-c-1$ positive integers
have sum at least $n-c-1$, whence $a\le n-(n-c-1)=c+1$.
\end{proof}

\begin{theorem}[Sharp uniform conductor range]\label{thm:sharp-conductor}
Let $D\subset\Pj^1$ have length $d\ge2$, let $1\le k<d$, and suppose
\begin{equation}\label{eq:sharp-range} n\ge2d-k.\end{equation}
Let $A\subset H^0(\mathcal O_D(n))$ be a generating $k$-plane, and put
$W=H^0(\mathcal O(n)(-D))$ and $U=\operatorname{res}_D^{-1}(A)$.
For every $m\ge2$,
\begin{equation}\label{eq:sharp-saturation}
 WU^{m-1}=H^0(\mathcal O(mn)(-D)).
\end{equation}
On the entire relative generating Grassmannian over
$\operatorname{Hilb}^d(\Pj^1)$, restriction gives
\begin{equation}\label{eq:sharp-cokernel}
 \operatorname{coker}(\Sym^m\mathcal U\to V_{mn}\otimes\mathcal O)
 \simeq
 \operatorname{coker}(\Sym^m\mathcal A\to\pi_*\mathcal O_{\mathcal D}(mn)).
\end{equation}
These are isomorphisms of coherent sheaves after arbitrary base change;
all their Fitting ideals agree. For a particular geometric fibre, the
weaker condition $n-d\ge\max_i a_i(U)-1$ suffices for the same conclusion.
Write $X_{d,k}^\circ$ for this relative generating Grassmannian. The bound \eqref{eq:sharp-range} is sharp as a bound uniform in $D,A$
and simultaneously in all $m\ge2$. For $k=d$ the cokernel assertion is immediate
from $U=V_n$ in the range $n\ge d-1$; the saturation assertion
in that case holds for $n\ge d$.
\end{theorem}
\begin{proof}
Write $D=(g)$, so $W=gV_{n-d}$. The series $U$ is base-point-free:
$W$ generates away from $D$ and $A$ generates on $D$.
Its codimension in $V_n$ is $c=d-k$. Lemma~\ref{lem:sharp-kernel}
and $n-d\ge c$ give
$WU=gV_{2n-d}$. Further multiplication by $U$ uses input degrees
$2n-d,3n-d,\ldots$, all at least $c$, and proves
\eqref{eq:sharp-saturation}. The same argument uses $a-1$ in place of
$c$ on a particular fibre.

All restriction targets and kernels are vector bundles: their first
cohomology groups vanish in this range. The map
$\mathcal W\otimes\Sym^{m-1}\mathcal U\to\pi_*I_{\mathcal D}(mn)$
is onto on every geometric fibre, hence is a surjection of vector
bundles. It remains onto after every base change. Locally split
$0\to\mathcal W\to\mathcal U\to\mathcal A\to0$. The kernel of the
map on symmetric powers is the image of
$\mathcal W\otimes\Sym^{m-1}\mathcal U$. Quotienting the restriction
square by these kernels proves \eqref{eq:sharp-cokernel}. Fitting
base change now applies \cite{StacksFitting}.

For sharpness take $n=2d-k-1$, $D=d\{0\}$, and, in the affine
coordinate $t$, take the degree-$n$ homogenizations of
\[
 U=\langle1,t,\ldots,t^{k-1},t^d,t^{d+1},\ldots,t^n\rangle.
\]
This is base-point-free and is the inverse image of a generating
$k$-plane on $D$. Exponents of its quadratic products lie in
\[
 [0,2k-2]\ \cup\ [d,n+k-1]\ \cup\ [2d,2n].
\]
None is $2d-1$. Thus $t^{2d-1}\notin U^2$, although it belongs to
$H^0(I_D(2n))$. Its nonzero cokernel class restricts to zero on $D$,
so even the cokernel comparison fails at this boundary in degree two.
This does not assert a sharp threshold separately for each $m>2$.
\end{proof}

The numerical bound can be replaced by a cohomological hypothesis on
a curve. Here and below an ideal such as $I_D$ is the ideal on the
curve, whereas a Fitting ideal is on the parameter space.

\begin{proposition}[An exact cohomological test]\label{prop:curve-kernel-test}
Let $C$ be a smooth projective curve, let $L$ be a line bundle, and
let $U\subset H^0(L)$ be base-point-free. Put $W=H^0(L(-D))$.
If $H^1(L(-D))=0$, then
\[
 \operatorname{coker}(W\otimes U\to H^0(L^2(-D)))
       \simeq H^1(M_U\otimes L(-D)).
\]
More generally, for $B_j=U^j\subset H^0(L^j)$ and its evaluation
kernel $M_{B_j}$, the obstruction to
$WB_j=H^0(L^{j+1}(-D))$ is exactly
$H^1(M_{B_j}\otimes L(-D))$.
\end{proposition}
\begin{proof}
Tensor the evaluation sequence for $B_j$ by $L(-D)$ and take
cohomology. The next term is $B_j\otimes H^1(L(-D))=0$.
\end{proof}

\begin{theorem}[Conductor transport on curves]\label{thm:curve-conductor}
Let $C$ have genus $g$, and let $D$ be effective of length $d\ge2$.
Suppose $L(-D)$ is globally generated,
$\Sym^2H^0(L(-D))\to H^0(L^2(-2D))$ is surjective, and
\begin{equation}\label{eq:curve-vanishing}
 H^1(L^j(-2D))=0\quad(j\ge1).
\end{equation}
For every generating $A\subset H^0(L|_D)$ and every $m\ge2$,
$U=\operatorname{res}_D^{-1}(A)$ satisfies
\[
 WU^{m-1}=H^0(L^m(-D)),\qquad
 \operatorname{coker}(\Sym^mU\to H^0(L^m))
 \simeq\operatorname{coker}(\Sym^mA\to H^0(L^m|_D)).
\]
In particular these assertions hold if
\begin{equation}\label{eq:curve-numerical}
 \deg L\ge2d+2g-1.
\end{equation}
They hold relatively, and after arbitrary base change, for smooth
projective families of curves with these fibrewise hypotheses and
locally free direct images of the displayed line bundles.
\end{theorem}
\begin{proof}
The hypotheses imply $H^1(L^j(-D))=0$ by the exact sequence with
quotient on $D$. They make $U$ base-point-free. It contains a pencil
without base points, obtained by avoiding the finitely many zeros of
one section. The pencil sequence, tensored by $L^j(-2D)$, shows that
multiplication onto $H^0(L^{j+1}(-2D))$ is surjective for $j\ge2$:
its obstruction is $H^1(L^{j-1}(-2D))$. Starting with the assumed
surjectivity for $W^2$, this proves
$W^2U^{m-2}=H^0(L^m(-2D))$.

The map $W\to H^0(L(-D)|_D)$ is onto by \eqref{eq:curve-vanishing}.
Choose $a\in A$ a unit on $D$ and lift it to $u\in U$.
Multiplication by $u^{m-1}$ identifies the first normal quotients in
these degrees. Consequently $Wu^{m-1}$ fills that quotient in
$H^0(L^m(-D))$. Together with the double-vanishing space just proved,
this gives the assertion about $WU^{m-1}$. The locally split
restriction square gives the cokernel isomorphism as before.

Under \eqref{eq:curve-numerical}, the line bundle $L(-D)$ has degree
at least $2g+1$. The classical normal-generation theorem in this
degree range gives the required quadratic surjectivity and global
generation; see \cite[p.~73, Introduction]{GL86}.
Serre duality gives \eqref{eq:curve-vanishing}, since
$\deg L(-2D)\ge2g-1$. The relative assertion follows from fibrewise
surjectivity of maps of vector bundles, exactly as in
Theorem~\ref{thm:sharp-conductor}.
\end{proof}

For example on any elliptic curve, a line bundle of degree at least
$2d+1$ transports all the finite-contact schemes below. The curve
statement is not a sharpness assertion in positive genus. Its
intrinsic hypotheses, and the exact obstruction in
Proposition~\ref{prop:curve-kernel-test}, distinguish a failure of
regularity from a failure of the finite-contact multiplication.
